\section{引言}
{\color{red}作者不对模板能够通过学校/学院的格式审查做任何明示或暗示的保证\footnote{这是一个脚注，脚注内容为小五号宋体，1.5倍行距。}。} 因参考文献使用GB7714标准，文档中还包含了Zeping Lee（李泽平） 编写的参考文献的样式文件（即两个后缀为.bst的文件）。

\section{系统要求}
宏包通过 CTeX 宏包来获得中文支持。 CTeX 宏包提供了一个
统一的中文\LaTeX 文档框架, 它可以从\url{http://www.ctex.org}网站下载。

此外， 宏包还使用了宏包 amsmath、 amsthm、 amsfonts、
amssymb、 bm 和 hyperref。

\subsection{下载与安装}
宏包的最新版本可以从\url{https://github.com/sslchi/lixinthesis}下载。

宏包包含一个文档类型文件.cls以及一个使用模板文件main.tex（也即本文的源文件）, 用户
可以通过修改这个模板来编写自己的学位论文。


\subsection{问题反馈}
因为毕业论文的要求很详细，从字体到到段距、行距等，细节不知凡几。因此难免有疏漏之处{\color{red}（再次申明，作者不对模板能够通过学校/学院的格式审查做任何明示或暗示的保证）}，因此用户在使用中如果发现问题，可以联系:
\begin{center}
公众号：MathICU
\end{center}

{\large \color{red}如果有同学愿意接手这个模板的维护，那是再好不过的了。}
\subsubsection{三级标题}
使用\verb*|\subssection|可以得到一个带编号的三级标题。

\paragraph{四级标题}
使用\verb*|\paragraph|可以得到一个带编号的三级标题。


\section{基本选项及命令介绍}\label{sec:basic}
本章介绍模板的一些基本选项和命令。关于\LaTeX 的基本用法，数学写作应该已经学过。
\subsection{编译}
本模板基于cetxart宏包，暂时只支持xeLaTeX方式进行编译。使用{\bfseries TexStudio编辑器}的用户可以直接点击编辑器上的构建并查看按钮进行编译，无需进行设置； 使用其它编辑器的用户请设置编译方式为xeLaTeX.

目前，只在TeXLive 2023上测试过，如果不能运行，请先检查版本是否是TeXLive 2023。

此外， 本文档中的编码方式皆为UTF-8,若用编辑器打开之后显示乱码，请将编辑器编码方式设置为UTF-8。


\subsection{模板中的选项及命令}
\begin{tabular}{p{10cm}p{5cm}}
	\verb|\title|         & 标题\\
	\verb|\englishtitle|  & 英文标题\\
	\verb|\IDnumber|      & 学号\\
	\verb|\author|        & 作者\\
	\verb|\school|        & 学院\\
	\verb|\major|         & 专业\\
	\verb|\class|         & 班级\\ 
	\verb|\submityear|    & 提交年\\
	\verb|\submitmonth|   & 提交月\\
\end{tabular}

